\section{Around Krull-Akizuki}
\label{section-krull-akizuki}

\noindent
One application of Krull-Akizuki is to show that there are plenty
of discrete valuation rings. More generally in this section we
show how to construct discrete valuation rings dominating Noetherian
local rings.

\medskip\noindent
First we show how to dominate a Noetherian local domain
by a $1$-dimensional Noetherian local domain by blowing up
the maximal ideal.

\begin{lemma}
\label{lemma-dominate-by-dimension-1}
Let $R$ be a local Noetherian domain with fraction field $K$.
Assume $R$ is not a field.
Then there exist $R \subset R' \subset K$ with
\begin{enumerate}
\item $R'$ local Noetherian of dimension $1$,
\item $R \to R'$ a local ring map, i.e., $R'$ dominates $R$, and
\item $R \to R'$ essentially of finite type.
\end{enumerate}
\end{lemma}

\begin{proof}
Choose any valuation ring $A \subset K$ dominating $R$ (which
exist by Lemma \ref{lemma-dominate}).
Denote $v$ the corresponding valuation.
Let $x_1, \ldots, x_r$ be a minimal set of generators of the
maximal ideal $\mathfrak m$ of $R$. We may and do assume that
$v(x_r) = \min\{v(x_1), \ldots, v(x_r)\}$. Consider the ring
$$
S = R[x_1/x_r, x_2/x_r, \ldots, x_{r - 1}/x_r] \subset K.
$$
Note that $\mathfrak mS = x_rS$ is a principal ideal.
Note that $S \subset A$ and that $v(x_r) > 0$, hence we see
that $x_rS \not = S$. Choose a minimal prime $\mathfrak q$
over $x_rS$. Then $\text{height}(\mathfrak q) = 1$ by
Lemma \ref{lemma-minimal-over-1}
and $\mathfrak q$ lies over $\mathfrak m$. Hence we
see that $R' = S_{\mathfrak q}$ is a solution.
\end{proof}

\begin{lemma}[Koll\'ar]
\label{lemma-hart-serre-loc-thm}
\begin{reference}
This is taken from a forthcoming paper by
J\'anos Koll\'ar entitled ``Variants of normality for Noetherian schemes''.
\end{reference}
Let $(R, \mathfrak m)$ be a local Noetherian ring.
Then exactly one of the following holds:
\begin{enumerate}
\item $(R, \mathfrak m)$ is Artinian,
\item $(R, \mathfrak m)$ is regular of dimension $1$,
\item $\text{depth}(R) \geq 2$, or
\item there exists a finite ring map $R \to R'$ which is not
an isomorphism whose kernel and cokernel are annihilated by a power
of $\mathfrak m$ such that $\mathfrak m$ is not an associated
prime of $R'$ and $R' \not = 0$.
\end{enumerate}
\end{lemma}

\begin{proof}
Observe that $(R, \mathfrak m)$ is not Artinian if and only if
$V(\mathfrak m) \subset \Spec(R)$ is nowhere dense. See
Proposition \ref{proposition-dimension-zero-ring}. We assume this from now on.

\medskip\noindent
Let $J \subset R$ be the largest ideal killed by a power of $\mathfrak m$.
If $J \not = 0$ then $R \to R/J$ shows that $(R, \mathfrak m)$
is as in (4).

\medskip\noindent
Otherwise $J = 0$. In particular $\mathfrak m$ is not an associated prime
of $R$ and we see that there is a nonzerodivisor $x \in \mathfrak m$ by
Lemma \ref{lemma-ideal-nonzerodivisor}. If $\mathfrak m$
is not an associated prime of $R/xR$ then $\text{depth}(R) \geq 2$
by the same lemma. Thus we are left with the case when there is a
$y \in R$, $y \not \in xR$ such that $y \mathfrak m \subset xR$.

\medskip\noindent
If $y \mathfrak m \subset x \mathfrak m$ then we can consider the
map $\varphi : \mathfrak m \to \mathfrak m$, $f \mapsto yf/x$
(well defined as $x$ is a nonzerodivisor). By the determinantal trick
of Lemma \ref{lemma-charpoly-module} there exists a monic
polynomial $P$ with coefficients in $R$ such that $P(\varphi) = 0$.
We conclude that $P(y/x) = 0$ in $R_x$.
Let $R' \subset R_x$ be the ring generated by
$R$ and $y/x$. Then $R \subset R'$ and $R'/R$ is a finite $R$-module
annihilated by a power of $\mathfrak m$. Thus $R$ is as in (4).

\medskip\noindent
Otherwise there is a $t \in \mathfrak m$ such that $y t = u x$
for some unit $u$ of $R$. After replacing $t$ by $u^{-1}t$
we get $yt = x$. In particular $y$ is a nonzerodivisor.
For any $t' \in \mathfrak m$ we have $y t' = x s$ for some $s \in R$.
Thus $y (t' - s t ) = x s - x s = 0$.
Since $y$ is not a zero-divisor this implies that $t' = ts$ and so
$\mathfrak m = (t)$. Thus $(R, \mathfrak m)$ is regular of dimension 1.

\medskip\noindent
The argument so far shows that every $R$ falls into one of the 4 cases.
To finish we have to show that no two of the properties can hold
simultaneously for each of the 6 combinations of properties. We'll
just show that (2) and (3) each exclude (4); the other cases are left
to the reader.

\medskip\noindent
Assume $R$ is regular of dimension $1$ and that $R \to R'$ is a finite
ring map whose kernel and cokernel are annihilated by a power of $\mathfrak m$
and $\mathfrak m$ is not an associated prime of $R'$. Then $R'$
is a finite $R$-module of depth at least $1$ and hence free by
Lemma \ref{lemma-regular-mcm-free}. Since $R \to R'$ is an isomorphism
in the generic point, we conclude that $R'$ must be free of rank $1$
as an $R$-module. Then $R' \to \text{End}_R(R') \cong R$ is an inverse
to the map $R \to R'$. Thus (4) cannot be true.

\medskip\noindent
Assume $R$ has depth $\geq 2$ and that $R \to R'$ is a finite
ring map whose kernel and cokernel are annihilated by a power
of $\mathfrak m$ and $\mathfrak m$ is not an associated prime of $R'$.
Then $R \to R'$ is necessarily injective (as the kernel would have
both depth $0$ and depth $\geq 1$) and the cokernel has depth
at least $1$ (by Lemma \ref{lemma-depth-in-ses}
and the fact that $R'$ has depth $\geq 1$ by assumption)
whence cannot be supported on
$\{\mathfrak m\}$ unless it is $0$ as well. Thus (4) cannot be true.
\end{proof}

\begin{lemma}
\label{lemma-nonregular-dimension-one}
Let $R$ be a local ring with maximal ideal $\mathfrak m$.
Assume $R$ is Noetherian, has dimension $1$, and that
$\dim(\mathfrak m/\mathfrak m^2) > 1$. Then there exists
a ring map $R \to R'$ such that
\begin{enumerate}
\item $R \to R'$ is finite,
\item $R \to R'$ is not an isomorphism,
\item the kernel and cokernel of $R \to R'$ are annihilated by
a power of $\mathfrak m$, and
\item $\mathfrak m$ is not an associated prime of $R'$.
\end{enumerate}
\end{lemma}

\begin{proof}
This follows from Lemma \ref{lemma-hart-serre-loc-thm}
and the fact that $R$ is not Artinian, not regular, and
does not have depth $\geq 2$ (the last part because the
depth does not exceed the dimension by
Lemma \ref{lemma-bound-depth}).
\end{proof}

\begin{example}
\label{example-nonreduced}
Consider the Noetherian local ring
$$
R = k[[x, y]]/(y^2)
$$
It has dimension 1 and it is Cohen-Macaulay.
An example of an extension as in
Lemma \ref{lemma-nonregular-dimension-one}
is the extension
$$
k[[x, y]]/(y^2) \subset k[[x, z]]/(z^2), \ \ y \mapsto xz
$$
in other words it is gotten by adjoining $y/x$ to $R$. The
effect of repeating the construction $n > 1$ times is
to adjoin the element $y/x^n$.
\end{example}

\begin{example}
\label{example-bad-dvr-char-p}
Let $k$ be a field of characteristic $p > 0$ such that $k$
has infinite degree over its subfield $k^p$ of $p$th powers.
For example $k = \mathbf{F}_p(t_1, t_2, t_3, \ldots)$.
Consider the ring
$$
A =
\left\{
\sum a_i x^i \in k[[x]] \text{ such that }
[k^p(a_0, a_1, a_2, \ldots) : k^p] < \infty
\right\}
$$
Then $A$ is a discrete valuation ring and its completion is
$A^\wedge = k[[x]]$. Note that the induced extension of fraction
fields of $A \subset k[[x]]$ is infinite purely inseparable.
Choose any $f \in k[[x]]$, $f \not \in A$. Let $R = A[f] \subset k[[x]]$.
Then $R$ is a Noetherian local domain of dimension $1$ whose
completion $R^\wedge$ is nonreduced (think!).
\end{example}

\begin{remark}
\label{remark-resolution-dim-1}
Suppose that $R$ is a $1$-dimensional semi-local Noetherian domain.
If there is a maximal ideal $\mathfrak m \subset R$ such that
$R_{\mathfrak m}$ is not regular, then we may apply
Lemma \ref{lemma-nonregular-dimension-one} to $(R, \mathfrak m)$
to get a finite ring extension $R \subset R_1$.
(For example one can do this so that $\Spec(R_1) \to \Spec(R)$
is the blowup of $\Spec(R)$ in the ideal $\mathfrak m$.)
Of course $R_1$ is a $1$-dimensional semi-local Noetherian
domain with the same fraction field as $R$. If $R_1$ is not a
regular semi-local ring, then we may repeat the construction to
get $R_1 \subset R_2$. Thus we get a sequence
$$
R \subset R_1 \subset R_2 \subset R_3 \subset \ldots
$$
of finite ring extensions which may stop if $R_n$ is regular for
some $n$. Resolution of singularities would be the claim
that eventually $R_n$ is indeed regular. In reality this is not
the case. Namely, there exists a characteristic $0$
Noetherian local domain $A$ of dimension $1$ whose completion is nonreduced,
see \cite[Proposition 3.1]{Ferrand-Raynaud} or our
Examples, Section \ref{examples-section-local-completion-nonreduced}.
For an example in characteristic $p > 0$ see
Example \ref{example-bad-dvr-char-p}.
Since the construction of blowing up commutes with completion it
is easy to see the sequence never stabilizes.
See \cite{Bennett} for a discussion (mostly in positive characteristic).
On the other hand, if the completion of $R$ in all of its maximal
ideals is reduced, then the procedure stops (insert future reference
here).
\end{remark}

\begin{lemma}
\label{lemma-characterize-dvr}
Let $A$ be a ring. The following are equivalent.
\begin{enumerate}
\item The ring $A$ is a discrete valuation ring.
\item The ring $A$ is a valuation ring and Noetherian but not a field.
\item The ring $A$ is a regular local ring of dimension $1$.
\item The ring $A$ is a Noetherian local domain with maximal ideal
$\mathfrak m$ generated by a single nonzero element.
\item The ring $A$ is a Noetherian local normal domain of dimension $1$.
\end{enumerate}
In this case if $\pi$ is a generator of the maximal ideal of
$A$, then every nonzero element of $A$ can be uniquely written as
$u\pi^n$, where $u \in A$ is a unit.
\end{lemma}

\begin{proof}
The equivalence of (1) and (2) is
Lemma \ref{lemma-valuation-ring-Noetherian-discrete}.
Moreover, in the proof of Lemma \ref{lemma-valuation-ring-Noetherian-discrete}
we saw that if $A$ is a discrete valuation ring, then $A$ is a PID, hence (3).
Note that a regular local ring is a domain (see
Lemma \ref{lemma-regular-domain}). Using this the equivalence of (3) and (4)
follows from dimension theory, see Section \ref{section-dimension}.

\medskip\noindent
Assume (3) and let $\pi$ be a generator of the maximal ideal $\mathfrak m$.
For all $n \geq 0$ we have
$\dim_{A/\mathfrak m} \mathfrak m^n/\mathfrak m^{n + 1} = 1$
because it is generated by $\pi^n$ (and it cannot be zero).
In particular $\mathfrak m^n = (\pi^n)$ and
the graded ring $\bigoplus \mathfrak m^n/\mathfrak m^{n + 1}$
is isomorphic to the polynomial ring $A/\mathfrak m[T]$.
For $x \in A \setminus \{0\}$ define
$v(x) = \max\{n \mid x \in \mathfrak m^n\}$.
In other words $x = u \pi^{v(x)}$ with $u \in A^*$.
By the remarks above we have $v(xy) = v(x) + v(y)$
for all $x, y \in A \setminus \{0\}$. We extend this to the field of fractions
$K$ of $A$ by setting $v(a/b) = v(a) - v(b)$ (well defined by multiplicativity
shown above). Then it is clear that $A$ is the set of elements of
$K$ which have valuation $\geq 0$. Hence we see that $A$ is a
valuation ring by Lemma \ref{lemma-valuation-valuation-ring}.

\medskip\noindent
A valuation ring is a normal domain by Lemma \ref{lemma-valuation-ring-normal}.
Hence we see that the equivalent conditions (1) -- (3) imply
(5). Assume (5). Suppose that $\mathfrak m$ cannot be generated
by $1$ element to get a contradiction.
Then Lemma \ref{lemma-nonregular-dimension-one} implies there is a finite
ring map $A \to A'$ which is an isomorphism after inverting
any nonzero element of $\mathfrak m$ but not an isomorphism.
In particular we may identify $A'$ with a subset of the fraction field of $A$.
Since $A \to A'$ is finite it is integral (see
Lemma \ref{lemma-finite-is-integral}).
Since $A$ is normal we get $A = A'$ a contradiction.
\end{proof}

\begin{definition}
\label{definition-uniformizer}
Let $A$ be a discrete valuation ring. A {\it uniformizer} is an element
$\pi \in A$ which generates the maximal ideal of $A$.
\end{definition}

\noindent
By Lemma \ref{lemma-characterize-dvr} any two uniformizers of a discrete
valuation ring are associates.

\begin{lemma}
\label{lemma-finite-length}
Let $R$ be a domain with fraction field $K$.
Let $M$ be an $R$-submodule of $K^{\oplus r}$.
Assume $R$ is local Noetherian of dimension $1$.
For any nonzero $x \in R$ we have $\text{length}_R(R/xR) < \infty$
and
$$
\text{length}_R(M/xM) \leq r \cdot \text{length}_R(R/xR).
$$
\end{lemma}

\begin{proof}
If $x$ is a unit then the result is true. Hence we may assume
$x \in \mathfrak m$ the maximal ideal of $R$. Since $x$ is not
zero and $R$ is a domain we have $\dim(R/xR) = 0$, and hence
$R/xR$ has finite length. Consider $M \subset K^{\oplus r}$ as
in the lemma. We may assume that the elements of $M$ generate
$K^{\oplus r}$ as a $K$-vector space after replacing $K^{\oplus r}$
by a smaller subspace if necessary.

\medskip\noindent
Suppose first that $M$ is a finite $R$-module. In that case we can clear
denominators and assume $M \subset R^{\oplus r}$. Since
$M$ generates $K^{\oplus r}$ as a vectors space we see that
$R^{\oplus r}/M$ has finite length. In particular there exists
an integer $c \geq 0$ such that $x^cR^{\oplus r} \subset M$.
Note that $M \supset xM \supset x^2M \supset \ldots$ is a sequence of
modules with successive quotients each isomorphic to $M/xM$. Hence
we see that
$$
n \text{length}_R(M/xM) = \text{length}_R(M/x^nM).
$$
The same argument for $M = R^{\oplus r}$ shows that
$$
n \text{length}_R(R^{\oplus r}/xR^{\oplus r}) =
\text{length}_R(R^{\oplus r}/x^nR^{\oplus r}).
$$
By our choice of $c$ above we see that $x^nM$ is sandwiched between
$x^n R^{\oplus r}$ and $x^{n + c}R^{\oplus r}$. This easily gives that
$$
r(n + c) \text{length}_R(R/xR)
\geq
n \text{length}_R(M/xM)
\geq
r (n - c) \text{length}_R(R/xR)
$$
Hence in the finite case we actually get the result of the lemma with
equality.

\medskip\noindent
Suppose now that $M$ is not finite. Suppose that the length of $M/xM$ is
$\geq k$ for some natural number $k$. Then we can find
$$
0 \subset N_0 \subset N_1 \subset N_2 \subset \ldots N_k \subset M/xM
$$
with $N_i \not = N_{i + 1}$ for $i = 0, \ldots k - 1$.
Choose an element $m_i \in M$ whose congruence class mod $xM$ falls
into $N_i$ but not into $N_{i - 1}$ for $i = 1, \ldots, k$.
Consider the finite $R$-module $M' = Rm_1 + \ldots + Rm_k \subset M$.
Let $N'_i \subset M'/xM'$ be the inverse image of $N_i$.
It is clear that $N'_i \not =N'_{i + 1}$ by our choice of $m_i$.
Hence we see that $\text{length}_R(M'/xM') \geq k$. By the
finite case we conclude $k \leq r\text{length}_R(R/xR)$
as desired.
\end{proof}

\noindent
Here is a first application.

\begin{lemma}
\label{lemma-finite-extension-residue-fields-dimension-1}
Let $R \to S$ be a homomorphism of domains inducing an
injection of fraction fields $K \subset L$. If $R$ is Noetherian
local of dimension $1$ and $[L : K] < \infty$ then
\begin{enumerate}
\item each prime ideal $\mathfrak n_i$ of $S$ lying over
the maximal ideal $\mathfrak m$ of $R$ is maximal,
\item there are finitely many of these, and
\item $[\kappa(\mathfrak n_i) : \kappa(\mathfrak m)] < \infty$ for each $i$.
\end{enumerate}
\end{lemma}

\begin{proof}
Pick $x \in \mathfrak m$ nonzero. Apply Lemma \ref{lemma-finite-length}
to the submodule $S \subset L \cong K^{\oplus n}$ where $n = [L : K]$.
Thus the ring $S/xS$ has finite length over $R$. It follows that
$S/\mathfrak m S$ has finite length over $\kappa(\mathfrak m)$.
In other words, $\dim_{\kappa(\mathfrak m)} S/\mathfrak m S$
is finite (Lemma \ref{lemma-dimension-is-length}). Thus $S/\mathfrak mS$
is Artinian (Lemma \ref{lemma-finite-dimensional-algebra}). The
structural results on Artinian rings implies parts (1) and (2), see
for example Lemma \ref{lemma-artinian-finite-length}.
Part (3) is implied by the finiteness established above.
\end{proof}

\begin{lemma}
\label{lemma-finite-length-global}
Let $R$ be a domain with fraction field $K$.
Let $M$ be an $R$-submodule of $K^{\oplus r}$.
Assume $R$ is Noetherian of dimension $1$.
For any nonzero $x \in R$ we have
$\text{length}_R(M/xM) < \infty$.
\end{lemma}

\begin{proof}
Since $R$ has dimension $1$ we see that $x$ is contained in finitely many
primes $\mathfrak m_i$, $i = 1, \ldots, n$, each maximal. Since $R$ is
Noetherian we see that $R/xR$ is Artinian and
$R/xR = \prod_{i = 1, \ldots, n} (R/xR)_{\mathfrak m_i}$ by
Proposition \ref{proposition-dimension-zero-ring} and
Lemma \ref{lemma-artinian-finite-length}. Hence $M/xM$ similarly
decomposes as the product $M/xM = \prod (M/xM)_{\mathfrak m_i}$
of its localizations at the $\mathfrak m_i$. By Lemma \ref{lemma-finite-length}
applied to $M_{\mathfrak m_i}$ over $R_{\mathfrak m_i}$ we see each
$M_{\mathfrak m_i}/xM_{\mathfrak m_i} = (M/xM)_{\mathfrak m_i}$
has finite length over $R_{\mathfrak m_i}$. Thus $M/xM$
has finite length over $R$ as the above implies $M/xM$
has a finite filtration by $R$-submodules whose successive
quotients are isomorphic to the residue fields $\kappa(\mathfrak m_i)$.
\end{proof}

\begin{lemma}[Krull-Akizuki]
\label{lemma-krull-akizuki}
Let $R$ be a domain with fraction field $K$.
Let $L/K$ be a finite extension of fields.
Assume $R$ is Noetherian and $\dim(R) = 1$.
In this case any ring $A$ with $R \subset A \subset L$ is
Noetherian.
\end{lemma}

\begin{proof}
Let $I \subset A$ be a nonzero ideal. By
Lemma \ref{lemma-intersection-not-zero}
we can find a nonzero element $x \in I \cap R$.
Then we get $I/xA \subset A/xA$. By Lemma \ref{lemma-finite-length-global}
the $R$-module $A/xA$ has finite length as an $R$-module. Hence
$I/xA$ has finite length as an $R$-module. Hence $I$ is finitely
generated as an ideal in $A$.
\end{proof}

\begin{lemma}
\label{lemma-exists-dvr}
Let $R$ be a Noetherian local domain with fraction field $K$.
Assume that $R$ is not a field.
Let $L/K$ be a finitely generated field extension.
Then there exists discrete valuation ring $A$ with fraction field
$L$ which dominates $R$.
\end{lemma}

\begin{proof}
If $L$ is not finite over $K$ choose a transcendence basis
$x_1, \ldots, x_r$ of $L$ over $K$ and replace $R$ by
$R[x_1, \ldots, x_r]$ localized at the maximal ideal
generated by $\mathfrak m_R$ and $x_1, \ldots, x_r$.
Thus we may assume $K \subset L$ finite.

\medskip\noindent
By Lemma \ref{lemma-dominate-by-dimension-1} we may assume $\dim(R) = 1$.

\medskip\noindent
Let $A \subset L$ be the integral closure of $R$ in $L$.
By Lemma \ref{lemma-krull-akizuki} this is Noetherian.
By Lemma \ref{lemma-integral-overring-surjective} there
is a prime ideal $\mathfrak q \subset A$ lying
over the maximal ideal of $R$.
By Lemma \ref{lemma-characterize-dvr} the ring $A_{\mathfrak q}$ is a discrete
valuation ring dominating $R$ as desired.
\end{proof}








